\chapter{Almost-everywhere varifold blow-ups}\label{sec:rectifiable-varifold-tangents}
A rectifiable varifold carries both a weighted measure and a tangent-plane field. Scalar convergence of the weight does not by itself control the plane coordinate. We prove the needed approximate continuity and obtain the full varifold blow-up at almost every point. The conventions are those of \cite{Sim18}, Chapter 3, pp.~90--92, and Chapter 8, Section 1, pp.~235--236. In this section $1\leq m<n$, $M\subset U\subset\mathbb R^n$ is measurable and countably $m$-rectifiable with locally finite unweighted area, and $0<\theta<\infty$ is measurable with $\mu=\theta\mathcal H^m\llcorner M$ locally finite.

\begin{lemma}[A Borel representative of the tangent-plane field]\label{lem:rectifiable-tangent-field-borel}
There is a Borel matrix field $A:U\to\mathbb R^{n\times n}$ which equals the orthogonal projection onto $T_xM$ at $\mathcal H^m$-almost every $x\in M$. It can be chosen to be a rank-$m$ orthogonal projection at every point of $U$. Any two such representatives agree $\mu$-almost everywhere on $M$.
\end{lemma}
\begin{proof}
Use the Borel representative of $M$ and the countable $C^1$ patch cover from Theorem \ref{thm:rectifiable-c1-decomposition} and Lemma \ref{lem:finite-area-c1-patches}. Assign each point to the first patch containing it. On a patch with regular derivative matrix $DF$, projection onto its tangent is
\[
 DF\bigl((DF)^*DF\bigr)^{-1}(DF)^*.
\]
This is continuous in chart coordinates; the inverse chart is continuous, so it is a Borel field on the patch. The first-patch assignment gives a Borel field on the covered part of $M$. Put any fixed rank-$m$ projection on its complement. By Theorem \ref{thm:rectifiable-forward-tangents}, the assigned field is the approximate tangent almost everywhere. Uniqueness of that tangent gives agreement of any two representatives there; $\mu\ll\mathcal H^m\llcorner M$ gives weighted almost-everywhere agreement. No differentiability of the tangent field across different pieces is asserted.
\end{proof}

\begin{proposition}[The tangent field is approximately continuous for the weight]\label{prop:rectifiable-weighted-tangent-oscillation}
For the representative in Lemma \ref{lem:rectifiable-tangent-field-borel}, at $\mu$-almost every $x\in M$,
\[
 \lim_{r\downarrow0}\frac1{\mu(B_r(x))}
       \int_{M\cap B_r(x)}\|A(y)-A(x)\|_{\mathrm{op}}\,d\mu(y)=0.
\]
At the same almost-everywhere points the normalized tilt integral satisfies
\[
 \lim_{r\downarrow0}r^{-m}\int_{M\cap B_r(x)}
       \|A(y)-A(x)\|_{\mathrm{op}}^2\,d\mu(y)=0.
\]
\end{proposition}
\begin{proof}
Each of the finitely many matrix entries of $A$ is bounded and locally $\mu$-integrable. The absolute-oscillation version of Theorem \ref{thm:radon-function-differentiation} applies to all of them on one full-measure set. The operator norm is bounded by the sum of the absolute matrix entries, which gives the first limit. The difference of two orthogonal projections has norm at most one: for a unit vector $v$, the quadratic form of the symmetric matrix $A-B$ lies between $-1$ and $1$, and the spectral theorem gives this operator norm bound. Thus the squared norm is at most the norm. Theorem \ref{thm:weighted-rectifiable-forward} gives $\mu(B_r)=O(r^m)$ at almost every point, and multiplication yields the second limit. Positivity of $\theta$ makes the weighted and unweighted exceptional sets equivalent.
\end{proof}

\begin{theorem}[A rectifiable varifold has a unique planar blow-up almost everywhere]\label{thm:rectifiable-varifold-ae-blowup}
Let $V=V(M,\theta)$ in the sense of Definition \ref{def:rectifiable-varifold}. For $\mu$-almost every $x\in M$, its dilations satisfy
\[
 (\eta_{x,r})_\#V\rightharpoonup\theta(x)|T_xM|
 \quad\text{as }r\downarrow0,
\]
where $|P|$ is the multiplicity-one varifold of the linear plane $P$. In particular the convergence holds at every vanishing scale, as a measure on position and plane coordinates. No stationarity or bound on first variation is needed for this almost-everywhere assertion.
\end{theorem}
\begin{proof}
Choose $x$ where Proposition \ref{prop:weighted-tangent-local-area}, the density limit and Proposition \ref{prop:rectifiable-weighted-tangent-oscillation} hold, and where $A(x)$ represents $P=T_xM$. Let $\Phi(z,Q)$ be a continuous Grassmann-bundle test whose position support lies in $B_R(0)$. The dilation formula from Definition \ref{def:varifold-dilation} gives
\[
 \int\Phi\,d((\eta_{x,r})_\#V)
 =r^{-m}\int_M\Phi((y-x)/r,T_yM)\,d\mu(y).
\]
The planes are unchanged by scalar dilation, and its tangential Jacobian is $r^{-m}$. Compare this with the same integral with $P$ in the plane argument. On the compact position ball times the compact Grassmannian, $\Phi$ is uniformly continuous. For any $\delta>0$, the part where $d(T_yM,P)\leq\delta$ contributes a difference at most its uniform modulus of continuity at $\delta$ times $r^{-m}\mu(B_{Rr}(x))$. This mass factor stays bounded. On the remaining part, Markov's inequality bounds the difference by
\[
 \frac{2\|\Phi\|_\infty}{\delta}\,r^{-m}
 \int_{B_{Rr}(x)}\|A(y)-A(x)\|_{\mathrm{op}}\,d\mu(y)\longrightarrow0.
\]
First let $r\downarrow0$ and then $\delta\downarrow0$. The plane-coordinate error vanishes. Proposition \ref{prop:weighted-tangent-local-area}, applied to the scalar continuous test $z\mapsto\Phi(z,P)$, gives the remaining limit $\theta(x)\int_P\Phi(z,P)\,d\mathcal H^m(z)$. This is exactly the asserted plane varifold. The full small-radius limits in the two propositions give full convergence, not merely subsequential compactness.
\end{proof}

\begin{corollary}[Integral multiplicity equals almost-everywhere density]\label{cor:integral-varifold-density-multiplicity}
If $V=V(M,\theta)$ is integral, then for $\|V\|$-almost every $x$,
\[
 \Theta^m(\|V\|,x)=\theta(x)\in\{1,2,3,\ldots\},
 \qquad (\eta_{x,r})_\#V\rightharpoonup\theta(x)|T_xM|.
\]
This is an almost-everywhere multiplicity statement; it does not assert that every support point has integer density or a unique tangent cone.
\end{corollary}
\begin{proof}
The weighted density limit follows from Proposition \ref{prop:weighted-tangent-local-area} by the unit-ball cutoff argument. Theorem \ref{thm:rectifiable-varifold-ae-blowup} gives the plane-coordinate convergence. Integrality is exactly the positive integer condition on $\theta$ almost everywhere. Since $\theta>0$ and is locally integrable, Hausdorff-null and weighted-null subsets of $M$ agree. Thus the conclusions hold at almost every point for the weight. Singular intersections and other exceptional points are not covered by this conclusion.
\end{proof}
